\documentclass[12pt]{article}

\usepackage{lipsum}
\usepackage[margin=1in]{geometry}
\usepackage{graphicx}
\graphicspath{{../data}}
\usepackage{float} % for H specifier
\usepackage{subcaption} % subfigures

\usepackage{xcolor}
\begin{document}
%\twocolumn
\section{Multiple images}
\lipsum[2]

\begin{figure}[!h]%[h] % ignoring text ratio...
  \centering
  % ---- row 1 ----
  \begin{subfigure}[b]{0.48\linewidth}
    \centering
    \includegraphics[width=\linewidth]{images/blue.jpg}
    \caption{Top left}\label{fig:blueCross}
  \end{subfigure}%
  \hfill
  %\hspace{10mm}
  \begin{subfigure}[b]{0.48\linewidth}
    \centering
    \includegraphics[width=\linewidth]{images/cyan.jpg}
    \caption{Top right}\label{fig:cyanCross}
  \end{subfigure}

  \vspace{6pt}    % vertical gap between rows

  % ---- row 2 ----
  \begin{subfigure}[b]{0.48\linewidth}
    \centering
    \includegraphics[width=\linewidth]{images/green.jpg}
    \caption{Bottom left}\label{fig:greenCross}
  \end{subfigure}%
  \hfill
  \begin{subfigure}[b]{0.48\linewidth}
    \centering
    \includegraphics[width=\linewidth]{images/orange.jpg}
    \caption{Bottom right}\label{fig:orangeCross}
  \end{subfigure}
  \caption{A 2$\times$2 grid of images.}
  \label{fig:grid}
\end{figure}

\textbf{
The blue panel is in the top left of Figure \ref{fig:tabGrid} it is the Figure \ref{fig:blueCross}
}

\lipsum[3]

\begin{figure}[!h]
    \centering
    \begin{tabular}{c c}
        \includegraphics[width=0.28\linewidth]{images/blue.jpg} &  \includegraphics[width=0.28\linewidth]{images/cyan.jpg}\\
        (a) Top left & (b) Top  right\\
        \includegraphics[width=0.28\linewidth]{images/green.jpg} &  \includegraphics[width=0.28\linewidth]{images/orange.jpg}\\
        (c) Bottom left & (d) Bottom right
    \end{tabular}
    \caption{A $2\times2$ table of images.}\label{}
    \label{fig:tabGrid}
\end{figure}
\textbf{The blue panel is in the top left of Figure \ref{fig:tabGrid}}


\lipsum[5-7]

%\onecolumn

\section{Picture}
\lipsum[1-5]

\begin{figure}[!h]
    \centering
       \begin{picture}(15cm, 13cm)%10cm for overlap
      % BLUE
      \put(0,10mm){\includegraphics[width=0.28\linewidth]{images/blue.jpg}}
      \put(16mm, 5mm){\parbox{12mm}{\centering \footnotesize{Caption below} }}
    
      % CYAN........
      \put(90mm, 10mm){\includegraphics[width=0.28\linewidth]{images/cyan.jpg}}
      \put(96mm,0.35\linewidth){Cyan caption above}
    
      % RED
      \put(10mm,0.4\linewidth){\includegraphics[width=0.28\linewidth]{images/red.jpg}}
      \put(5mm,80mm){\rotatebox{90}{On the side}}
      
      % ORANGE rotate  
      \put(75mm,65mm){\includegraphics[width=0.28\linewidth,angle=30]{images/orange.jpg}}
      \put(70mm,90mm){{\rotatebox{-60}{\parbox{25mm}{\footnotesize Inclined caption because why not}}}}
    \end{picture} 
    \caption{Figure arranged using the \texttt{picture} environment}
    \label{fig:picFigure}
\end{figure}

\lipsum[7-9]

\end{document}